\documentclass[14pt,a4paper]{extarticle}
% --- Настройки шрифтов и языка ---
\usepackage{fontspec}
\usepackage{polyglossia}
\setdefaultlanguage{russian}
\setotherlanguage{english}
% Подключаем Times New Roman
\setmainfont{Times New Roman}
\newfontfamily\cyrillicfont{Times New Roman}
\newfontfamily\cyrillicfonttt{Courier New} % для листингов кода
\newfontfamily\codefont{Courier New}
% --- Поля страницы (ГОСТ) ---
\usepackage[left=30mm, right=15mm, top=20mm, bottom=20mm]{geometry}
% --- Межстрочный интервал (1.5) ---
\usepackage{setspace}
\onehalfspacing
% --- Абзацный отступ (1.25 см) ---
\usepackage{indentfirst}
\setlength{\parindent}{1.25cm}
% --- Нумерация страниц ---
\usepackage{fancyhdr}
\pagestyle{fancy}
\fancyhf{}
\fancyfoot[C]{\thepage}
\renewcommand{\headrulewidth}{0pt}
% --- Настройка заголовков ---
\usepackage{titlesec}
\titleformat{\section}
  {\normalfont\fontsize{14pt}{14pt}\bfseries\raggedright\MakeUppercase}
  {\hspace*{1.25cm}\thesection}{1em}{}
\titleformat{\subsection}
  {\normalfont\fontsize{14pt}{14pt}\bfseries\raggedright}
  {\hspace*{1.25cm}\thesubsection}{1em}{}
\titleformat{\subsubsection}
  {\normalfont\fontsize{14pt}{14pt}\bfseries\raggedright}
  {\hspace*{1.25cm}\thesubsubsection}{1em}{}
\titlespacing*{\section}{0pt}{18pt}{12pt}
\titlespacing*{\subsection}{0pt}{18pt}{12pt}
\titlespacing*{\subsubsection}{0pt}{18pt}{12pt}
% Команда для не нумеруемых структурных элементов (ВВЕДЕНИЕ и т.д.)
\newcommand{\structelement}[1]{
\clearpage
\begin{center}
\bfseries\MakeUppercase{#1}
\end{center}
\addcontentsline{toc}{section}{#1}
\vspace{6pt}
}
% --- Оформление рисунков и таблиц ---
\usepackage{graphicx}
\usepackage{caption}
\usepackage{float}
% Запрет автоматических переносов слов по слогам
\hyphenpenalty=10000
\exhyphenpenalty=10000
\sloppy
\DeclareCaptionLabelFormat{rusfigure}{Рисунок #2}
\DeclareCaptionLabelSeparator{dashsep}{~---~}
\captionsetup[figure]{labelformat=rusfigure, labelsep=dashsep, justification=centering, font=normalsize, position=bottom}
\DeclareCaptionLabelFormat{rustable}{Таблица #2}
\captionsetup[table]{labelformat=rustable, labelsep=dashsep, justification=raggedright, singlelinecheck=false, font=normalsize, position=top}
% --- Дополнительные пакеты ---
\usepackage{amsmath, amssymb} % Формулы
\usepackage{tabularx}         % Таблицы
\usepackage{listings}        % Листинги кода
\usepackage{ragged2e}
\usepackage{xcolor}          % Цвета для кода
\usepackage{comment}
\lstset{
  basicstyle=\small\codefont,
  breaklines=true,
  frame=single,
  numbers=left,
  numberstyle=\tiny,
  showstringspaces=false,
  keepspaces=true,
  columns=fullflexible
}
\setlength{\emergencystretch}{1em}
\usepackage{microtype}
% --- Оглавление ---
\usepackage{tocloft}
\renewcommand{\cftsecleader}{\cftdotfill{\cftdotsep}} % точки в оглавлении
\usepackage{hyperref}
\hypersetup{colorlinks=false, pdfborder={0 0 0}}

% =====================================================
% НАЧАЛО ДОКУМЕНТА
% =====================================================
\begin{document}

% Титульный лист
\thispagestyle{empty} % Убираем номер страницы и колонтитулы
\begin{center}
    \textbf{Министерство науки и высшего образования Российской Федерации}
    
    \vspace{5pt}

    ФЕДЕРАЛЬНОЕ ГОСУДАРСТВЕННОЕ АВТОНОМНОЕ ОБРАЗОВАТЕЛЬНОЕ УЧРЕЖДЕНИЕ ВЫСШЕГО ОБРАЗОВАНИЯ
    
    \vspace{5pt}
    
    \textbf{НАЦИОНАЛЬНЫЙ ИССЛЕДОВАТЕЛЬСКИЙ УНИВЕРСИТЕТ ИТМО}
    
    \vspace{5pt}
    
    \textbf{ITMO University}

    \vspace{15pt}

    \textbf{ИССЛЕДОВАТЕЛЬСКАЯ РАБОТА}
\end{center}

\begin{flushleft}
    \textbf{По теме:} Logistic map \\
    \textbf{По дисциплине:} Математический анализ и основы вычислений \\
    \textbf{Обучающийся:} Чепелкин Владислав Андреевич \\
    \textbf{Направление подготовки:} 11.03.02 Инфокоммуникационные технологии и системы связи \\
    \textbf{Образовательная программа:} Инженерия искусственного интеллекта
    
    \vspace{1cm}
    \textbf{Обучающийся}\qquad\rule{3.1cm}{0.01mm}\qquad\rule{2.5cm}{0.01mm}
    \quad\quad\underline{Чепелкин В. А.} \\
    \hspace{5.6cm}\textsuperscript{(дата)}
    \hspace{2.4cm}\textsuperscript{(подпись)}
    \hspace{2.1cm}\textsuperscript{(Ф.И.О.)}
    
    \textbf{Руководитель}\quad\quad\rule{3.1cm}{0.01mm}\quad\quad\rule{2.5cm}{0.01mm}
    \qquad\underline{Ершов А. Р.} \\
    \hspace{5.6cm}\textsuperscript{(дата)}
    \hspace{2.4cm}\textsuperscript{(подпись)}
    \hspace{2.1cm}\textsuperscript{(Ф.И.О.)}
\end{flushleft}

\vfill
\begin{center}
\fontsize{14pt}{14pt}\selectfont
Санкт-Петербург \\
2026
\end{center}

\clearpage

\begin{center}
    \bfseries СОДЕРЖАНИЕ
\end{center}
\vspace{12pt}
\renewcommand{\contentsname}{}
\vspace{-30pt}
\tableofcontents



\structelement{Цель работы}

Исследовать одномерные дискретные динамические системы на примере логистического отображения. Рассмотреть влияние управляющего параметра $r$ на вид отображения и сравнить логистическое отображение с отображением, соответствующим индивидуальному варианту.

Номер варианта:
\[
N = ISU\ mod\ 5 = 502527\ mod\ 5 = 2.
\]

Тогда
\[
x_{n+1} = g(x_n) = rx_n (1 - x_n)^2, \quad r \in \Bigl[ 0; \frac{27}{4} \Bigr].
\]

\clearpage


\section{Easy level}

\subsection{Сохранение интервала (0,1)}

Рассмотрим логистическое отображение
\[
x_{n+1} = r x_n(1 - x_n),
\]
где $0 < x_0 < 1$ и $r \in [0,1]$.
Требуется доказать, что
\[
0 < x_n < 1 \quad \forall n \in \mathbb{N}
\]

Сначала рассмотрим $r \in (0,1]$. Проведем доказательство методом математической индукции. По условию начальное значение удовлетворяет $0 < x_0 < 1$. Предположим, что для некоторого $n$ выполнено $0 < x_n < 1$. Отсюда $0 < 1 - x_n < 1$. Так как $r > 0$, имеем
\[
x_{n+1} = r x_n(1 - x_n) > 0.
\]
Для верхней границы воспользуемся известной оценкой
\[
x(1 - x) \leq \frac{1}{4}, \quad x \in [0,1].
\]
Ее можно непосредственно получить из равенства
\[
x(1 - x) = \frac{1}{4} - \Bigl( x - \frac{1}{2} \Bigr)^2 \leq \frac{1}{4}.
\]
Следовательно,
\[
x_{n+1} = r x_n(1 - x_n) \leq \frac{r}{4} \leq \frac{1}{4} < 1.
\]

Таким образом, $0 < x_{n+1} < 1$. Индукционный переход выполнен, поэтому при $r \in (0,1]$
\[
\forall n \geq 0:\ 0 < x_n < 1.
\]
Отдельно необходимо рассмотреть $r = 0$. В этом случае $x_1 = 0 \cdot x_0(1 - x_0) = 0$, и далее $x_2 = x_3 = ... = 0$. Следовательно, записанное в задании утверждение для $r = 0$ со строгим неравенством неверно. Для всего указанного диапазона $r \in [0,1]$ корректно утверждение
\[
0 \leq x_n < 1,
\]
а исходное строгое неравенство выполняется при $r \in (0,1]$.


\subsection{Влияние параметра $r$ на логистическое отображение}

Рассмотрим зависимость $x_n$ от $x_{n-1}$:
\[
x_n = r x_{n-1}(1 - x_{n-1}).
\]
Для удобства обозначим $f_r(x) = rx(1 - x) = rx - rx^2$. При $r > 0$ это парабола, ветви которой направлены вниз. Нули функции определяются из $rx(1 - x) = 0$, следовательно, $x = 0,\ x = 1$. Найдем положение максимума. Производная равна $f'_r(x) = r(1 - 2x)$. Условие $f'_r(x) = 0$ дает $x = 1/2$. Так как $f''_r(x) = -2r < 0$, в этой точке действительно достигается максимум. Его значение: $f_r(1/2) = r/4$.

Таким образом, при изменении $r$ положение нулей $x=0$ и $x=1$ и положение максимума $x=1/2$ остаются неизменными. Меняется высота параболы: $\max f_r(x) = r/4$. Чем больше $r$, тем сильнее график растягивается вдоль вертикальной оси.

Построим функции для нескольких значений $r$ из полного диапазона логистического отображения $r \in [0,4]$.

\begin{lstlisting}[language=python]
import numpy as np
import matplotlib.pyplot as plt

def logistic(x, r):
    return r * x * (1 - x)

x = np.linspace(0, 1, 500)

for r in [0.5, 1, 2, 3, 4]:
    plt.plot(x, logistic(x, r), label=f"r = {r}")

plt.xlabel(r"$x_{n-1}$")
plt.ylabel(r"$x_n$")
plt.title('Логистическое отображение')
plt.grid()
plt.legend()
plt.show()
\end{lstlisting}

\begin{figure}[H]
    \centering
    \includegraphics[scale=0.8]{Figure_1.png}
    \caption{Логистическое отображение}
\end{figure}

Параметр $r$ определяет вертикальный масштаб графика логистического отображения. При увеличении $r$ максимальное возможное значение $x_n$ увеличивается линейно и равно $r/4$. При этом общий вид параболы, ее нули и положение вершины не изменяются.


\subsection{Отображение функции $g(x_n) = r x_n (1 - x_n)^2,\ r \in [0, \frac{27}{4}]$}

Для $N=2$ задано отображение
\[
x_{n+1} = g(x_n) = r x_n(1 - x_n)^2, \quad r \in \Bigl[ 0, \frac{27}{4} \Bigr]
\]
Построим зависимость $x_n$ от $x_{n-1}$ для нескольких значений $r$.

\begin{lstlisting}[language=python]
def g(x, r):
    return r * x * (1 - x)**2

x = np.linspace(0, 1, 500)

for r in [1, 3, 5, 6.75]:
    plt.plot(x, g(x, r), label=f"r = {r}")

plt.xlabel(r"$x_{n-1}$")
plt.ylabel(r"$x_n$")
plt.title(r'Отображение $g(x)=rx(1-x)^2$')
plt.grid()
plt.legend()
plt.show()
\end{lstlisting}

\begin{figure}[H]
    \centering
    \includegraphics[scale=0.8]{Figure_2.png}
    \caption{Отображение $g(x) = rx(1-x)^2$}
\end{figure}

Исследуем форму полученных кривых. Функция
\[
g(x) = rx(1 - x)^2
\]
обращается в ноль в точках $x = 0,\ x = 1$. Определим положение максимума:
\[
g'(x) = r \bigl[ (1 - x)^2 - 2x(1 - x) \bigr] = r(1 - x)(1 - 3x).
\]
Критическими точками являются $x = 1/3,\ x = 1$. На интервале $(0, 1/3)$ производная положительна, а на $(1/3, 1)$ отрицательна. Следовательно, максимум на $[0,1]$ достигается при $x = 1/3$. Его значение:
\[
g \Bigl( \frac{1}{3} \Bigr) = r \cdot \frac{1}{3} \cdot \Bigl( 1 - \frac{1}{3} \Bigr)^2 = \frac{r}{3} \cdot \frac{4}{9} = \frac{4r}{27}.
\]
При максимально допустимом $r=27/4$ получаем
\[
g_{max} = \frac{4}{27} \cdot \frac{27}{4} = 1.
\]
Поэтому указанный в задании диапазон: $0 \leq r \leq 27/4$ обеспечивает $0 \leq g(x) \leq 1$ для всех $x \in [0,1]$. То есть отображение переводит отрезок $[0,1]$ в себя.

Параметр $r$, как и в логистическом отображении, масштабирует функцию по вертикальной оси: при увеличении $r$ ее максимальное значение линейно возрастает.


\subsection{Сравнение логистического отображения и отображения варианта}

Сравним
\[
f(x) = rx(1 - x)
\]
и
\[
g(x) = rx(1 - x)^2.
\]
У отображений есть несколько общих свойств:
\begin{itemize}
    \item обе функции непрерывны на $[0,1]$;
    \item обе неотрицательны на $[0,1]$;
    \item обе имеют нули при $x=0$ и $x=1$;
    \item обе сначала возрастают, достигают единственного внутреннего максимума, а затем убывают;
    \item в обоих случаях увеличение $r$ приводит к вертикальному растяжению графика.
\end{itemize}
При этом существуют и различия. Для логистического отображения максимум достигается при $x = 1/2$, а для отображения варианта – при $x = 1/3$. Максимальные значения также различаются: $f_{max} = r/4,\ g_{max} = 4r/27$.

Кроме того, поведение функций около $x = 1$ различается. Для логистического отображения множитель $(1 - x)$ имеет первую степень, поэтому при $r > 0$: $f'(1) = -r \neq 0$, то есть график пересекает ось $Ox$. У отображения варианта множитель $(1 - x)$ возведен в квадрат: $g'(1) = 0$. Поэтому при $x = 1$ график касается оси $Ox$, а не пересекает ее.

Для наглядного непосредственного сравнения при одном значении $r$ можно построить оба отображения на одном графике.

\begin{lstlisting}[language=python]
r = 3
x = np.linspace(0, 1, 500)

plt.plot(x, logistic(x, r), label=r"$rx(1-x)$")
plt.plot(x, g(x, r), label=r"$rx(1-x)^2$")

plt.xlabel(r"$x_{n-1}$")
plt.ylabel(r"$x_n$")
plt.title(f'Сравнение отображений при r = {r}')
plt.grid()
plt.legend()
plt.show()
\end{lstlisting}

\begin{figure}[H]
    \centering
    \includegraphics[scale=0.8]{Figure_3.png}
    \caption{Сравнение отображений при $r = 3$}
\end{figure}

Сходство отображений связано с их общей структурой: обе функции содержат множители $r$, $x$ и $(1 - x)$. Множитель $r$ задает вертикальный масштаб, а наличие $x$ и $(1 - x)$ обеспечивает нулевые значения на границах интервала.

Основное различие вызвано дополнительным множителем $(1 - x)$ в отображении варианта. Из-за него максимум смещается от $x = 1/2$ к $x = 1/3$, изменяется форма правой части графика и допустимый диапазон параметра $r$.

Таким образом, отображение варианта похоже на логистическое, однако его нелинейность имеет другую форму. Это позволяет предположить, что дальнейшее поведение двух динамических систем также будет иметь общие закономерности, но характерные значения управляющего параметра $r$ могут различаться.

\clearpage



\section{Normal level}

На данном уровне исследуем неподвижные точки, монотонность и сходимость последовательностей для логистического отображения, а затем проведем аналогичное исследование отображения варианта
\[
g(x) = rx(1 - x)^2.
\]


\subsection{Неподвижные точки логистического отображения}

Неподвижной называется точка $x^*$, для которой
\[
x^* = f(x^*).
\]
Для логистического отображения $f(x) = rx(1 - x)$ получаем уравнение $x^* = rx^*(1 - x^*)$. Перенесем $x^*$ в правую часть и вынесем общий множитель: $x^* \bigl[ r(1 - x^*) - 1 \bigr] = 0$. Отсюда $x^*_1 = 0$ или $r(1 - x^*) = 1$. При $r \neq 0$ второе уравнение дает $x^*_2 = 1 - \frac{1}{r} = \frac{r-1}{r}$.

Таким образом, неподвижные точки имеют вид
\[
x_1 = 0, \quad x_2 = 1 - \frac{1}{r}.
\]
При этом необходимо учитывать, что логистическое отображение рассматривается при $r \in [0,4]$ и $x \in [0,1]$.

При $0 \leq r < 1$ вторая точка либо не определена при $r = 0$, либо отрицательна при $0 < r < 1$: $1 - 1/r < 0$. Следовательно, в $[0,1]$ имеется только одна неподвижная точка $x^* = 0$.

При $r = 1$ обе ветви решения совпадают: $x^*_2 = 1 - 1 = 0$. То есть различная неподвижная точка по-прежнему одна.

При $r > 1$ имеем $0 < 1 - 1/r < 1$, поэтому на $[0,1]$ существуют две различные неподвижные точки: $x_1 = 0,\ x_2 = 1 - 1/r$.

Максимально логистическое отображение может иметь две различные неподвижные точки. Действительно, уравнение $rx(1 - x) = x$ после раскрытия скобок является квадратным: $-r x^2 + (r - 1)x = 0$, а квадратное уравнение не может иметь более двух различных действительных корней.


\subsection{Монотонность и предел при $r \in (0,1]$}

Пусть
\[
x_{n+1} = rx_n(1 - x_n), \quad x_0 \in (0,1),\ r \in (0,1].
\]
В Easy level было доказано, что в этом случае
\[
0 < x_n < 1 \quad \forall n.
\]
Проверим отношение двух соседних элементов:
\[
\frac{x_{n+1}}{x_n} = r(1 - x_n).
\]
Так как $0 < x_n < 1$, то $0 < 1 - x_n < 1$. С учетом $0 < r \leq 1$ получаем $0 < r(1 - x_n) < 1$. Следовательно, $x_{n+1} < x_n$.

Таким образом, последовательность $\{x_n\}$ строго монотонно убывает. При этом $x_n > 0$, то есть последовательность ограничена снизу нулем. По теореме о сходимости монотонной ограниченной последовательности существует конечный предел
\[
\lim_{n \rightarrow \infty} x_n = L, \quad L \geq 0.
\]
Найдем его. Переходя к пределу в рекуррентном соотношении $x_{n+1} = rx_n(1 - x_n)$, получаем благодаря непрерывности правой части $L = rL(1 - L)$. Следовательно, $L \bigl[ r(1 - L)-1 \bigr] = 0$.

Возможны два решения: $L = 0$ и, при $r \neq 0$, $L = 1 - 1/r$. Но при $0 < r \leq 1\ $ $\ 1 - 1/r \leq 0$. При $r < 1$ это значение отрицательно и не может быть пределом положительной последовательности. При $r = 1$ оно совпадает с нулем. Поэтому во всем рассматриваемом диапазоне
\[
\lim_{n \rightarrow \infty} x_n = 0.
\]

Итак, при $x_0 \in (0,1)$ и $r \in (0,1]$ последовательность $\{x_n\}$ строго монотонно убывает и сходится к нулю. Проверим результат графически.

\begin{lstlisting}[language=python]
import numpy as np
import matplotlib.pyplot as plt

def logistic(x, r):
    return r * x * (1 - x)

def trajectory(f, x0, r, steps=30):
    xs = [x0]
    for _ in range(steps):
        xs.append(f(xs[-1], r))
    return np.array(xs)

x0 = 0.8
n = np.arange(31)

for r in [0.3, 0.6, 0.9, 1.0]:
    xs = trajectory(logistic, x0, r, 30)
    plt.plot(n, xs, marker='o', markersize=3, label=f'r = {r}')

plt.xlabel('n')
plt.ylabel(r'$x_n$')
plt.title(r'Сходимость $x_n$ при $r \in (0,1]$')
plt.grid()
plt.legend()
plt.show()
\end{lstlisting}

\begin{figure}[H]
    \centering
    \includegraphics[scale=0.8]{Figure_4.png}
    \caption{Сходимость $x_n$ при $r \in (0,1]$}
\end{figure}

График подтверждает аналитический результат: последовательности монотонно убывают и стремятся к нулю.


\subsection{Четная и нечетная подпоследовательности при $r \in (2,3)$}

Пусть теперь
\[
r \in (2,3).
\]
Ненулевая неподвижная точка логистического отображения равна
\[
x^* = 1 - \frac{1}{r}.
\]
По условию задания предполагается, что
\[
x_{2n} > x^*,\ x_{2n+1} < x^*.
\]
Необходимо определить монотонность подпоследовательностей $\{x_{2n}\}$ и $\{x_{2n+1}\}$. Для этого рассмотрим отображение, выполняемое два раза: $F(x) = f \bigl( f(x) \bigr)$. Так как
\[
f(x) = rx(1 - x),\ F(x) = r \bigl[ rx(1 - x) \bigr] \bigl[ 1 - rx(1 - x) \bigr] = r^2 x(1 - x)[1 - rx(1 - x)].
\]
Рассмотрим разность $F(x)-x$. После разложения на множители получаем
\[
F(x) - x = -x(rx - r + 1) (r^2 x^2 - r^2 x - rx + r + 1).
\]
Первый множитель $rx - r + 1$ обращается в ноль при $x = (r - 1) / r = x^*$. Корни последнего квадратного множителя:
\[
x_{\pm} = \frac{r+1 \pm \sqrt{ \bigl( (r - 3)(r + 1) \bigr)}}{2r}.
\]
При $2 < r < 3$ выражение $(r - 3)(r + 1) < 0$, поэтому дополнительных действительных корней нет. Более того, квадратный множитель $Q(x) = r^2 x^2 - r(r + 1)x + r + 1$ положителен для всех действительных $x$, поскольку его старший коэффициент положителен, а дискриминант отрицателен.

Следовательно, при $x > 0$ знак $F(x) - x$ определяется множителем $-(rx - r + 1)$. Если $x > x^*$, то $rx - r + 1 > 0$, поэтому $F(x) - x < 0$. То есть $F(x) < x$.

Применяя это к $x = x_{2n} > x^*$, получаем
\[
x_{2n+2} < x_{2n}.
\]
Следовательно, четная подпоследовательность $\{x_{2n}\}$ убывает.

Если же $0 < x < x^*$, то $rx - r + 1 < 0$, откуда $F(x) - x > 0$ и поэтому $F(x) > x$. Для $x = x_{2n+1} < x^*$ получаем $x_{2n+3} > x_{2n+1}$. Следовательно, нечетная подпоследовательность $\{x_{2n+1}\}$ возрастает.

При выполнении условия задания две подпоследовательности приближаются к неподвижной точке с разных сторон:
\[
x_{2n} \downarrow x^*, \quad x_{2n+1} \uparrow x^*.
\]
Последнее согласуется с данным в условии фактом, что при $r \in (1,3]$ $\lim_{n \rightarrow \infty} x_n = x^*$.

Проверим результат графически, взяв, например, $r = 2,8$ и начальную точку выше $x^*$.

\begin{lstlisting}[language=python]
r = 2.8
x_star = 1 - 1 / r
x0 = 0.8

xs = trajectory(logistic, x0, r, 30)
n = np.arange(len(xs))

plt.plot(n[::2], xs[::2], 'o-', label=r'$x_{2n}$')
plt.plot(n[1::2], xs[1::2], 'o-', label=r'$x_{2n+1}$')
plt.axhline(x_star, color='black', linestyle='--', label=r'$x^*$')

plt.xlabel('n')
plt.ylabel(r'$x_n$')
plt.title(f'Подпоследовательности при r = {r}')
plt.grid()
plt.legend()
plt.show()
\end{lstlisting}

\begin{figure}[H]
    \centering
    \includegraphics[scale=0.8]{Figure_5.png}
    \caption{Подпоследовательности при $r = 2,8$}
\end{figure}

На графике четные элементы убывают к $x^*$, а нечетные возрастают к $x^*$. Полная последовательность при этом сходится к неподвижной точке немонотонно, совершая затухающие колебания вокруг нее.


\subsection{Неподвижные точки отображения варианта}

Для $N = 2$ имеем
\[
g(x) = rx(1 - x)^2.
\]
Неподвижные точки удовлетворяют
\[
x^* = r x^*(1 - x^*)^2.
\]
Вынесем $x^*$:
\[
x^* \bigl[ r(1 - x^*)^2 - 1 \bigr] = 0.
\]
Первая неподвижная точка: $x^*_0 = 0$. Для ненулевой точки необходимо решить $r(1 - x^*)^2 = 1$. При $r > 0$:
\[
(1 - x^*)^2 = \frac{1}{r},
\]
откуда алгебраически
\[
x^* = 1 \pm \frac{1}{\sqrt{r}}.
\]
Однако в рассматриваемой динамической системе $x \in [0,1]$. Точка $1 + 1/\sqrt{r} > 1$ лежит вне этого интервала. Поэтому возможная ненулевая неподвижная точка на $[0,1]$ имеет вид $x^*_1 = 1 - 1/\sqrt{r}$. Она принадлежит $[0,1]$ только при $r \geq 1$.

Таким образом:
\begin{itemize}
    \item при $0 \leq r \leq 1$ на $[0,1]$ имеется только неподвижная точка $x^* = 0$ (при $r = 1$ вторая формула также дает $0$);
    \item при $r > 1$ имеются две различные неподвижные точки: $0$ и $1 - 1 / \sqrt{r}$.
\end{itemize}
Для заданного диапазона $r \in [0,27/4]$ это полностью описывает неподвижные точки отображения.


\subsection{Диапазон монотонной сходимости к нулю для отображения варианта}

Рассмотрим
\[
x_{n+1} = rx_n(1 - x_n)^2, \quad x_0 \in (0,1)
\]
Для строгой монотонной сходимости к нулю требуется
\[
0 < x_{n+1} < x_n.
\]
Так как $x_n > 0$, разделим неравенство $x_{n+1} < x_n$ на $x_n$:
\[
r(1 - x_n)^2 < 1.
\]
Если $0 < r \leq 1$, то для любого $x_n \in (0,1)$
\[
0 < (1 - x_n)^2 < 1,
\]
поэтому
\[
0 < r(1 - x_n)^2 < 1.
\]
Следовательно,
\[
0 < x_{n+1} < x_n.
\]
Кроме того, $x_{n+1} = rx_n(1 - x_n)^2 \leq 4r/27 \leq 4/27 < 1$, поэтому все последующие элементы также остаются внутри $(0,1)$. Значит, при $0 < r \leq 1$ последовательность строго убывает и ограничена снизу нулем, а потому имеет предел $L \geq 0$.

Из непрерывности $g$:
\[
L = rL(1 - L)^2.
\]
Если $L \neq 0$, то
\[
r(1 - L)^2 = 1.
\]

При $r < 1$ такое равенство невозможно для $L \in [0,1]$. При $r = 1$ оно требует $L = 0$. Следовательно,
\[
\lim_{n \rightarrow \infty} x_n = 0.
\]

Покажем также, почему это является точной общей границей, если утверждение должно выполняться для любого $x_0 \in (0,1)$. При $r > 1$ выберем достаточно малое положительное $x_0$. Тогда
\[
r(1 - x_0)^2 > 1,
\]
и поэтому
\[
x_1 > x_0.
\]
Следовательно, при $r > 1$ нельзя гарантировать монотонное убывание к нулю для всех $x_0 \in (0,1)$.

Итак, общий диапазон параметра, в котором для любого начального $x_0 \in (0,1)$ последовательность монотонно сходится к нулю:
\[
0 < r \leq 1.
\]
При $r = 0$ последовательность после первой итерации также становится равной нулю:
\[
x_1 = 0,
\]
но переход $x_0 \rightarrow x_1$ нарушает строгую положительность. Если под монотонностью понимать нестрогое убывание, можно включить и $r = 0$.


\subsection{Графики зависимости $x_n$ от $n$ для отображения варианта}

В задании требуется построить графики для нескольких различных значений $r$. Возьмем значения как до, так и после $r = 1$, чтобы увидеть качественное изменение динамики.

\begin{lstlisting}[language=python]
def g(x, r):
    return r * x * (1 - x)**2

x0 = 0.2
steps = 40
n = np.arange(steps + 1)

for r in [0.5, 1.0, 2.0, 4.0, 6.0]:
    xs = trajectory(g, x0, r, steps)
    plt.plot(n, xs, marker='o', markersize=3, label=f'r = {r}')

plt.xlabel('n')
plt.ylabel(r'$x_n$')
plt.title(r'Траектории для $g(x)=rx(1-x)^2$')
plt.grid()
plt.legend()
plt.show()
\end{lstlisting}

\begin{figure}[H]
    \centering
    \includegraphics[scale=0.8]{Figure_6.png}
    \caption{Траектории для $g(x) = rx (1-x)^2$}
\end{figure}

При $r \leq 1$ наблюдается монотонное стремление последовательности к нулю, что соответствует полученному аналитическому результату.

При $r > 1$ появляется положительная неподвижная точка
\[
x^* = 1 - \frac{1}{\sqrt{r}},
\]
и характер динамики меняется: последовательность уже не обязана сходиться к нулю. При дальнейшем увеличении $r$ ненулевая неподвижная точка впоследствии теряет устойчивость, и характер динамики усложняется.

Для логистического отображения при $0 < r \leq 1$ последовательность монотонно убывает к нулю, а при $2 < r < 3$ сходимость к ненулевой неподвижной точке может происходить с чередованием относительно $x^*$: четная и нечетная подпоследовательности сходятся к ней с разных сторон. Для отображения варианта $g(x) = rx(1 - x)^2$ диапазон гарантированной монотонной сходимости к нулю для любого $x_0 \in (0,1)$ также равен $0 < r \leq 1$, а при $r > 1$ возникает ненулевая неподвижная точка $x^* = 1 - 1 / \sqrt{r}$.

\clearpage



\section{Hard level}

На данном уровне исследуем циклы логистического отображения, их изменение при росте параметра $r$, построим лестницу Ламерея, а затем проведем аналогичное исследование отображения варианта
\[
x_{n+1} = g(x_n) = r x_n(1 - x_n)^2.
\]


\subsection{Изменение длины цикла логистического отображения при $r \in (3, r_{\infty})$}

Рассмотрим логистическое отображение
\[
x_{n+1} = rx_n(1 - x_n)
\]
и интервал
\[
3 < r < r_{\infty}, \quad r_{\infty} \approx 3,5699456.
\]
При $r = 3$ неподвижная точка
\[
x^* = 1 - \frac{1}{r}
\]
теряет устойчивость, и возникает устойчивый цикл периода $2$. При дальнейшем увеличении $r$ этот цикл также теряет устойчивость и сменяется циклом периода $4$. Затем аналогично возникают циклы периодов $8$, $16$ и т. д.

Такой процесс называется каскадом удвоения периода. Первые значения $r$, при которых происходит смена периода, приближенно равны
\begin{itemize}
    \item $r_1 = 3$,
    \item $r_2 \approx 3,44949$,
    \item $r_3 \approx 3,54409$,
    \item $r_4 \approx 3,56441$
\end{itemize}
Все эти значения стремятся к $r_{\infty}$. Следовательно, при движении от $3$ к $r_{\infty}$ длина устойчивого цикла последовательно изменяется по закону
\[
m = 2,\ 4,\ 8,\ 16,\ 32,\ ...
\]

Экспериментально это можно проверить численно. Для определения периода сначала отбросим переходный участок траектории, после чего проверим, через сколько итераций значения начинают повторяться с заданной точностью.

\begin{lstlisting}[language=python]
import numpy as np
import matplotlib.pyplot as plt

def logistic(x, r):
    return r * x * (1 - x)

def find_period(f, r, x0=0.2, transient=5000,
                check=500, max_period=128, tol=1e-8):
    x = x0

    # Отбрасываем переходный процесс
    for _ in range(transient):
        x = f(x, r)

    # Сохраняем установившуюся часть траектории
    values = []
    for _ in range(check + max_period):
        x = f(x, r)
        values.append(x)

    values = np.array(values)

    # Ищем минимальный период
    for p in range(1, max_period + 1):
        if np.allclose(values[p:], values[:-p], atol=tol, rtol=0):
            return p

    return None
    
for r in [3.2, 3.5, 3.55, 3.565]:
    print(r, find_period(logistic, r))
\end{lstlisting}

\begin{table}[H]
\centering
\caption{Численный эксперимент каскада удвоения периода}
\begin{tabular}{|c|c|}
\hline
$r$ & Цикл периода \\
\hline
3.2   & 2  \\
3.5   & 4  \\
3.55  & 8  \\
3.565 & 16 \\
\hline
\end{tabular}
\end{table}

Конкретные значения около точек разветвления лучше выбирать не слишком близко к границам интервалов: непосредственно в их окрестности переходный процесс становится очень длинным, поэтому простое численное определение периода становится менее надежным.

Таким образом, численный эксперимент подтверждает каскад удвоения периода.


\subsection{Ограничения на порядок цикла $m$}

Из предыдущего исследования следует закономерность:
\[
m = 2^k,\ k = 1, 2, 3, ...
\]
То есть в рассматриваемом каскаде при $r \in (3, r_{\infty})$ устойчивые циклы имеют периоды, являющиеся степенями двойки: 2, 4, 8, 16, 32, ...

Это не означает, что уравнение $f^m(x) = x$ вообще не имеет других математических решений. Речь идет именно об установившемся устойчивом режиме, наблюдаемом при итерациях логистического отображения на основном каскаде удвоения периода.

По мере приближения $r$ к $r_{\infty}$ число удвоений неограниченно возрастает:
\[
m \rightarrow \infty\ \text{при}\ r \rightarrow r_{\infty} − 0.
\]
Точки бифуркаций накапливаются при $r \rightarrow r_{\infty}^-$; за $r_{\infty}$ начинается хаотическая область, содержащая также окна периодичности.

Для наглядности построим график установившейся части последовательности для четырех значений $r$.

\begin{lstlisting}[language=python]
def get_trajectory(f, x0, r, transient=1000, steps=40):
    x = x0

    for _ in range(transient):
        x = f(x, r)

    values = []
    for _ in range(steps):
        x = f(x, r)
        values.append(x)

    return np.array(values)

fig, axes = plt.subplots(2, 2, figsize=(10, 7))

r_values = [3.2, 3.5, 3.55, 3.565]
axes_flat = axes.flat

for i in range(len(r_values)):
    r = r_values[i]
    ax = axes_flat[i]

    values = get_trajectory(logistic, 0.2, r)

    ax.plot(values, 'o-', markersize=3)
    ax.set_title(f'r = {r}')
    ax.set_xlabel('n')
    ax.set_ylabel(r'$x_n$')
    ax.grid()

plt.tight_layout()
plt.show()
\end{lstlisting}

\begin{figure}[H]
    \centering
    \includegraphics[scale=0.65]{Figure_7.png}
    \caption{Графики для четырех значений $r$}
\end{figure}

На графиках последовательно наблюдаются 2, 4, 8 и 16 повторяющихся значений.


\subsection{Лестница Ламерея}

Лестница Ламерея позволяет наглядно представить итерационный процесс
\[
x_{n+1} = f(x_n).
\]
На одном графике строятся функция $y = f(x)$ и диагональ $y = x$. Начав с точки $(x_0, 0)$, перемещаемся вертикально до $(x_0, x_1)$, где $x_1 = f(x_0)$. Затем горизонтально переходим к диагонали: $(x_1, x_1)$. После этого снова движемся вертикально до графика функции: $(x_1, x_2)$, и повторяем процесс.

Напишем универсальную функцию, пригодную как для логистического отображения, так и для отображения варианта.

\begin{lstlisting}[language=python]
def cobweb(f, r, x0, steps=30, title='Лестница Ламерея'):
    x = np.linspace(0, 1, 1000)

    plt.figure(figsize=(7, 7))
    plt.plot(x, f(x, r), label=r'$f(x)$')
    plt.plot(x, x, 'k--', label=r'$y=x$')

    xn = x0
    yn = f(xn, r)

    # Первый шаг из (x0, 0)
    plt.plot([xn, xn], [0, yn], color='red')

    for _ in range(steps):
        # Из (xn, x_{n+1}) горизонтально к диагонали
        plt.plot([xn, yn], [yn, yn], color='red')

        xn = yn
        yn = f(xn, r)

        # Вертикально к графику отображения
        plt.plot([xn, xn], [xn, yn], color='red')

    plt.xlim(0, 1)
    plt.ylim(0, 1)
    plt.xlabel(r'$x_n$')
    plt.ylabel(r'$x_{n+1}$')
    plt.title(title)
    plt.grid()
    plt.legend()
    plt.show()

params = [
    (3.2, 20),
    (3.5, 40),
    (3.55, 60)
]
for r, steps in params:
    cobweb(
        logistic,
        r=r,
        x0=0.2,
        steps=steps,
        title=f'Лестница Ламерея: r = {r}'
    )
\end{lstlisting}

\begin{figure}[H]
    \centering
    \includegraphics[scale=0.4]{Figure_8.png}
    \caption{Лестница Ламерея с $r=3,2$}
\end{figure}

\begin{figure}[H]
    \centering
    \includegraphics[scale=0.4]{Figure_9.png}
    \caption{Лестница Ламерея с $r=3,5$}
\end{figure}

\begin{figure}[H]
    \centering
    \includegraphics[scale=0.4]{Figure_10.png}
    \caption{Лестница Ламерея с $r=3,55$}
\end{figure}


\subsection{Отображение циклов на лестнице Ламерея}

Для устойчивой неподвижной точки ступени лестницы постепенно приближаются к одной точке пересечения графиков
\[
y = f(x)\ \text{и}\ y = x.
\]
Для цикла периода 2 после переходного процесса лестница начинает повторять замкнутый маршрут между двумя различными значениями $x_1^*$ и $x_2^*$. На рисунке возникает характерный четырехугольный контур.

При периоде 4 траектория последовательно проходит через четыре различных значения
\[
x_1^* \rightarrow x_2^* \rightarrow x_3^* \rightarrow x_4^* \rightarrow x_1^*.
\]
Лестница соответственно образует более сложный повторяющийся замкнутый маршрут.

В общем случае для цикла порядка $m$ после исчезновения переходного процесса значения повторяются:
\[
x_1^* \rightarrow x_2^* \rightarrow ... \rightarrow x_m^* \rightarrow x_1^*.
\]
Поэтому на лестнице Ламерея наблюдается повторяющаяся замкнутая последовательность вертикальных и горизонтальных переходов через $m$ различных значений состояния. Чем больше период, тем сложнее визуально становится такой маршрут.

Таким образом, лестница Ламерея позволяет не только увидеть сходимость, но и визуально различать установившиеся циклы разных порядков.


\subsection{Исследование циклов отображения варианта}

Теперь рассмотрим
\[
g(x) = rx(1 - x)^2, \quad r \in \Bigl[ 0, \frac{27}{4} \Bigr]
\]
В Normal level была найдена ненулевая неподвижная точка
\[
x^* = 1 - \frac{1}{\sqrt{r}},
\]
которая существует на $[0,1]$ при $r > 1$.

Перед численным экспериментом полезно определить, при каком $r$ эта неподвижная точка теряет устойчивость. Это позволит понять, где следует искать первый цикл периода 2.

Производная:
\[
g'(x) = r(1 - x)(1 - 3x).
\]
В неподвижной точке
\[
x^* = 1 - \frac{1}{\sqrt{r}}
\]
имеем
\[
1 - x^* = \frac{1}{\sqrt{r}}.
\]
Кроме того,
\[
1 - 3x^* = 1 - 3 \Bigl( 1 - \frac{1}{\sqrt{r}} \Bigr) = -2 + \frac{3}{\sqrt{r}}.
\]
Следовательно,
\[
g'(x^*) = r \cdot \frac{1}{\sqrt{r}} \cdot (-2 + \frac{3}{\sqrt{r}}) = 3 - \frac{2}{\sqrt{r}}.
\]
Неподвижная точка устойчива при
\[
|g'(x^*)| < 1.
\]
То есть
\[
|3 - 2\sqrt{r}|<1.
\]
Получаем
\[
-1 < 3 - 2\sqrt{r} < 1.
\]
Из правой части:
\[
3 - 2\sqrt{r} < 1 \Rightarrow \sqrt{r} > 1 \Rightarrow r > 1.
\]
Из левой:
\[
3 - 2\sqrt{r} > -1 \Rightarrow \sqrt{r} < 2 \Rightarrow r < 4.
\]
Следовательно, положительная неподвижная точка устойчива при
\[
1 < r < 4.
\]
При $r = 4$ имеем
\[
g'(x^*) = -1,
\]
и происходит первая бифуркация удвоения периода. Поэтому при $r > 4$ ожидается появление устойчивого цикла периода 2.


\subsection{Численное исследование периода для $g(x)$}

Используем написанную ранее функцию $find\_period$.

\begin{lstlisting}[language=python]
def g(x, r):
    return r * x * (1 - x)**2

for r in [3.5, 4.2, 5.0, 5.5, 6.0, 6.5]:
    period = find_period(
        g, r,
        x0=0.2,
        transient=10000,
        check=1000,
        max_period=256,
        tol=1e-8
    )
    print(r, period)
\end{lstlisting}
Для систематического исследования просканируем весь диапазон параметра после первой бифуркации.

\begin{lstlisting}[language=python]
r_values = np.linspace(4.01, 6.74, 700)

periods = []

for r in r_values:
    p = find_period(
        g, r,
        x0=0.2,
        transient=5000,
        check=500,
        max_period=128,
        tol=1e-7
    )
    periods.append(p if p is not None else np.nan)
\end{lstlisting}
Поскольку периоды удваиваются, удобнее дополнительно использовать логарифмическую шкалу:

\begin{lstlisting}[language=python]
plt.figure(figsize=(9, 5))
plt.scatter(r_values, periods, s=5)

plt.yscale('log', base=2)
plt.xlabel('r')
plt.ylabel('Период m')
plt.title(r'Изменение периода цикла для $g(x)=rx(1-x)^2$')
plt.grid()
plt.show()
\end{lstlisting}

Здесь важно учитывать численный характер определения периода. В хаотических областях либо для периода выше $max\_period$ функция возвращает $None$. Кроме того, непосредственно около бифуркаций требуется особенно большое количество переходных итераций.

\begin{figure}[H]
    \centering
    \includegraphics[scale=0.65]{Figure_11.png}
    \caption{Изменения периода цикла для $g(x) = rx(1-x)^2$}
\end{figure}

Численный эксперимент показывает, что после $r = 4$ возникает цикл периода 2. Около $r = 5$ происходит следующее удвоение периода, после чего появляются циклы периодов 4, 8 и более высоких порядков. При дальнейшем росте $r$ наблюдаются также более сложные режимы, поэтому простой каскад удвоений уже не описывает весь диапазон параметра.


\subsection{График установившихся значений для отображения варианта}

Чтобы изменение режима было видно без погрешности алгоритма распознавания периода, построим также график последних значений траектории в зависимости от $r$.

Для каждого $r$ сначала выполним большое число итераций, чтобы отбросить переходный процесс, а затем нанесем последние значения $x_n$.

\begin{lstlisting}[language=python]
r_values = np.linspace(0.01, 6.75, 3000)

rs = []
xs = []

for r in r_values:
    x = 0.2

    for _ in range(2000):
        x = g(x, r)

    for _ in range(100):
        x = g(x, r)
        rs.append(r)
        xs.append(x)

plt.figure(figsize=(10, 6))
plt.plot(rs, xs, ',k', alpha=0.5)

plt.xlabel('r')
plt.ylabel(r'$x_n$')
plt.title(r'Установившиеся значения для $g(x)=rx(1-x)^2$')
plt.xlim(0, 6.75)
plt.ylim(0, 1)
plt.show()
\end{lstlisting}

\begin{figure}[H]
    \centering
    \includegraphics[scale=0.65]{Figure_12.png}
    \caption{Установившиеся значения для $g(x) = rx(1-x)^2$}
\end{figure}

Этот график нужен здесь для наглядного отображения изменения количества установившихся значений и, следовательно, длины цикла.


\subsection{Лестница Ламерея для отображения варианта}

Написанная функция $cobweb$ универсальна, поэтому дополнительных функций создавать не требуется. Например, построим лестницы до и после потери устойчивости ненулевой неподвижной точки:

\begin{lstlisting}[language=python]
for r in [3.5, 4.2, 5.0]:
    cobweb(
        g,
        r=r,
        x0=0.2,
        steps=60,
        title=rf'$g(x)=rx(1-x)^2$, r = {r}'
    )
\end{lstlisting}

\begin{figure}[H]
    \centering
    \includegraphics[scale=0.4]{Figure_13.png}
    \caption{$g(x) = rx(1-x)^2,\ r=3,5$}
\end{figure}

\begin{figure}[H]
    \centering
    \includegraphics[scale=0.4]{Figure_14.png}
    \caption{$g(x) = rx(1-x)^2,\ r=4,2$}
\end{figure}

\begin{figure}[H]
    \centering
    \includegraphics[scale=0.4]{Figure_15.png}
    \caption{$g(x) = rx(1-x)^2,\ r=5,0$}
\end{figure}

При $r < 4$ лестница стремится к ненулевой неподвижной точке. После $r = 4$ возникает цикл периода 2. При дальнейшем увеличении $r$ наблюдаются новые удвоения периода.


\subsection{Сравнение с логистическим отображением}

Численное исследование показывает качественное сходство двух отображений. Как для
\[
f(x) = rx(1 - x),
\]
так и для
\[
g(x) = rx(1 - x)^2
\]
при увеличении управляющего параметра устойчивая ненулевая неподвижная точка в некоторый момент теряет устойчивость, после чего возникает цикл периода 2. Затем наблюдается дальнейшее удвоение периода:
\[
2 \rightarrow 4 \rightarrow 8 \rightarrow 16 \rightarrow ...
\]

Основное различие состоит в значениях параметра $r$, при которых происходят эти изменения. Для логистического отображения первая бифуркация происходит при $r = 3$, а для отображения варианта – при $r = 4$.

Это различие связано с разной формой функций и, соответственно, с разными значениями их производных в неподвижных точках.

При этом обе функции на $[0,1]$ являются гладкими одномодальными отображениями: имеют один внутренний максимум. Поэтому сходство перехода от устойчивой неподвижной точки к циклам все более высоких периодов закономерно.

\clearpage



\structelement{Вывод}

В ходе работы было исследовано поведение одномерных дискретных динамических систем на примере логистического отображения
\[
x_{n+1} = rx_n(1 - x_n)
\]
и отображения индивидуального варианта $N = 2$
\[
x_{n+1} = rx_n(1 - x_n)^2.
\]

Было установлено, что управляющий параметр $r$ определяет не только масштаб графика отображения, но и характер поведения порождаемой последовательности. Для логистического отображения при $0 < r \leq 1$ последовательность с начальным условием $x_0 \in (0,1)$ монотонно убывает и сходится к нулевой неподвижной точке. При увеличении $r$ возникает ненулевая неподвижная точка $x^* = 1 - 1/r$, а при $2 < r < 3$ приближение к ней может происходить с чередованием относительно $x^*$: четная и нечетная подпоследовательности сходятся к неподвижной точке с разных сторон.

При дальнейшем увеличении параметра неподвижная точка теряет устойчивость. Для логистического отображения это происходит при $r = 3$, после чего возникает устойчивый цикл периода 2. На интервале $3 < r < r_{\infty}$, где $r_{\infty} \approx 3.5699456$, наблюдается каскад удвоения периода:
\[
2 \rightarrow 4 \rightarrow 8 \rightarrow 16 \rightarrow ...
\]
Таким образом, периоды устойчивых циклов основного каскада имеют вид $m = 2^k$. Лестница Ламерея позволила наглядно проследить переход от сходимости к одной неподвижной точке к циклическому поведению различных порядков.

Для отображения индивидуального варианта были получены качественно схожие результаты. Его ненулевая неподвижная точка
\[
x^* = 1 - \frac{1}{\sqrt{r}}
\]
возникает при $r > 1$ и устойчива при $1 < r < 4$. При $r = 4$ она теряет устойчивость, после чего возникает цикл периода 2, а при дальнейшем увеличении $r$ наблюдается последовательное появление циклов более высоких порядков и каскад удвоения периода.

Несмотря на различия в форме функций и значениях параметра, при которых происходят изменения режима, оба исследованных отображения демонстрируют сходный сценарий развития динамики: сходимость к неподвижной точке сменяется периодическим поведением, после чего при увеличении управляющего параметра возникают циклы все более высоких порядков.

Таким образом, на примере двух отображений было показано, что даже простое детерминированное рекуррентное соотношение способно демонстрировать различные режимы поведения. Аналитическое исследование неподвижных точек и сходимости в сочетании с численными экспериментами, графиками траекторий и лестницами Ламерея позволило проследить изменение этих режимов при изменении параметра $r$.



\end{document}